\documentclass[11pt]{article}
\usepackage[a4paper,margin=25mm]{geometry}
\usepackage{amsmath,amssymb}
\usepackage[hidelinks]{hyperref}
\setlength{\emergencystretch}{3em}
\title{Conjecture 00000001385: an infinite Hilbert spectral-count plateau}
\author{lizaixi01, with OpenAI Codex assistance}
\date{23 October 2026}
\begin{document}
\maketitle
\section{A necessary jump clause fails}
The source defines $N_T(\lambda)$ as a spectral counting
\emph{function} and requires that after an infinite value, all
higher values equal $2^{\aleph_0}$. With the standard Hilbert
dimension of spectral subspaces, this necessary clause is false.
Take the genuine probability circle
\[
 X=\mathbb R/\mathbb Z,\qquad \mu=\text{normalized Haar measure},
 \qquad T(x)=x.
\]
Its actual Koopman operator on $E=L^2(X,\mu;\mathbb C)$ is
$Uf=f\circ T=f$. Its only eigenvalue is the real number $1$,
and the eigenspace is all of $E$. The genuine Fourier Hilbert
basis is indexed by $\mathbb Z$. Consequently
\[
 N_T(1)=\aleph_0=N_T(2)\ne 2^{\aleph_0}.
\]
Here $1$ and $2$ are real \emph{thresholds}; $2$ is not asserted
to be a second spectral point. Thus the infinite-value jump
clause, and hence the source's conjunction under this standard
function reading, is disproved. We do not need to determine its
other claims about alternation or spectral-measure atomicity.

\section{Actual transformations and cutoff subspaces}
Lean constructs Mathlib's actual \texttt{AddCircle 1}, its Haar
probability measure, and the actual \texttt{Lp} space with exponent
$2$. The map $T$ is proved measure preserving and involutive.
The operator $U$ is the actual
\texttt{Lp.compMeasurePreserving} construction and its
complex-linear refinement, not an assigned diagonal model.
The proof establishes $U=\mathrm{id}$, existence of eigenvalue
$1$, and the absence of any other eigenvalue.

General complex eigenvalues have no specified order in the
source. This example avoids inventing one: its spectrum consists
only of the real eigenvalue $1$. The formal real-eigenvalue cutoff
is the genuine closed subspace
\[
 S_a=\overline{\sum_{r\in\mathbb R,\ r\leq a}\ker(U-rI)}.
\]
For robustness, the package independently constructs the modulus
cutoff
\[
 M_a=\overline{\sum_{z\in\mathbb C,\ |z|\leq a}\ker(U-zI)}.
\]
For every real $a\geq1$, the actual eigenspace at $1$ is one of
the summands and already equals $E$, so $S_a=M_a=E$.
Both cutoff interpretations give the displayed counterexample.

\newpage
\section{Dimension is proved from genuine Hilbert bases}
Mathlib's actual Fourier Hilbert basis is used with all its
orthonormality and completeness proofs. The formal argument
also proves that \emph{every} Hilbert basis of $E$ has index
cardinality $\aleph_0$, so its dimension is independent of the
chosen Fourier enumeration. If $(b_i)_{i\in I}$ is another
Hilbert basis, then for each Fourier vector $e_n$ the support of
its square-summable coordinate vector in $(b_i)$ is countable.
The union of those supports over $n\in\mathbb Z$ is all of $I$:
otherwise a nonzero $b_i$ would be orthogonal to every Fourier
vector, contradicting Fourier completeness. Thus $I$ is
countable. It cannot be finite, because that would make $E$
finite dimensional despite the infinite linearly independent
Fourier family. Hence $|I|=\aleph_0$.

The package defines the Hilbert dimension of a subspace by the
infimum of the cardinalities admitting an actual Hilbert basis.
For the relevant full subspace it proves that this set is the
nonempty singleton $\{\aleph_0\}$, using an explicitly transported
Fourier basis for existence and the preceding uniqueness proof.
There is no empty-infimum fallback, no assigned value for $N_T$,
and no substitution of Hamel \texttt{Module.rank} for Hilbert
dimension. The two counting functions are obtained by applying
this dimension definition to the actual cutoff subspaces.
Cantor's strict inequality $\aleph_0<2^{\aleph_0}$ closes both
formal negations of the required jump property.

\section{Conventions, primary references and verification}
The Hilbert interpretation is standard in this setting.
S. New's Waterloo \emph{Chapter 2: Hilbert Spaces}, PDF/printed
page 10, Theorem 2.35 and Definition 2.36, proves basis-cardinality
uniqueness and explicitly uses Hilbert dimension for a Hilbert
space unless otherwise specified, distinguished from Hamel
dimension. B. Solomyak's \emph{Ergodic Theory} notes,
21 November 2019, PDF page 1, formulate the actual Koopman
operator on $L^2$ and orthogonal eigenfunction bases.
The directly read primary sources are
\href{https://www.math.uwaterloo.ca/~snew/PMATH453/Chap2HilbertSpaces.pdf}{New's notes}
and
\href{https://u.math.biu.ac.il/~solomyb/TEACH/20/ET/spectral.pdf}{Solomyak's notes}.

The scope is explicit: this disproves the necessary infinite-value
jump for the printed counting \emph{function}, with Hilbert
dimension and real or modulus thresholds. It does not claim to
refute a different Hamel-dimension interpretation, an unspecified
ordering on arbitrary complex spectra, or a rule sampled only at
distinct actual spectral points. The identity example has only
one such point. These interpretations are not silently conflated.

The package pins Lean 4.33.0 and Mathlib commit
\begin{center}\small\texttt{db584cd6d46c92f209a44c0f1c829460d327499d}\end{center}
and supplies complete LaTeX/PDF, the original bilingual source,
all named declaration axiom audits, a new root build and full
official \texttt{leanchecker --fresh Main} replay. No
\texttt{sorry}, \texttt{native\_decide} or additional axiom is used.
OpenAI Codex assisted, with separate supervising AI proof/source/
report/PDF review. This is not human certification; organizer
acceptance is pending.
\end{document}
